\documentclass[11pt]{article}
\usepackage{xspace,mathrsfs}
\usepackage{graphicx,rotating}
%\usepackage{graphicx,pstricks,rotating}

\usepackage{url}
%\usepackage{latexsym}
\usepackage{amscd,amsfonts,amsmath,amssymb}
\usepackage{color}
\usepackage{longtable}
\usepackage{multirow,setspace}
\usepackage{algorithm}
\usepackage{algpseudocode}
\renewcommand{\topfraction}{0.9}
\renewcommand{\textfraction}{0.1}

\usepackage{float}
\usepackage{booktabs}
\usepackage{caption}
\usepackage{fancyhdr}


\def\shownotes{0}   % set 1 for version with author notes
                    % set 0 for no notes

%%%%%%%%%%%%%%%%%%%%%%%%%%%%%%%%%%%%%%%%%%%%%%%%%%%%%%%%%%%%%%
%Some macros (you can ignore everything until "end of macros")
\topmargin 0pt \advance \topmargin by -\headheight \advance
\topmargin by -\headsep

\textheight 8.9in
\oddsidemargin 0pt \evensidemargin \oddsidemargin \marginparwidth
0.5in

\textwidth 6.5in

%%%%%%
\DeclareMathAlphabet{\mathsl}{OT1}{cmr}{m}{sl}
\DeclareMathAlphabet{\mathsc}{OT1}{cmr}{m}{sc}
\DeclareMathAlphabet{\mathslbf}{OT1}{cmr}{bx}{sl}

\newcommand{\procfont}[1]{\mathsc{#1}}
\newcommand{\tablefont}[1]{\mathsf{#1}}
\newcommand{\Initialize}{\procfont{Initialize}}
\newcommand{\Finalize}{\procfont{Finalize}}
\newcommand{\FinalizeCodeNoarg}{\underline{\textrm{proc} $\Finalize$}}
\newcommand{\Onquery}[2]{\underline{\textrm{proc} $#1(#2)$}}
\newcommand{\InitializeCode}{\underline{\textrm{proc} $\Initialize$}}
\newcommand{\FinalizeCode}[1]{\underline{\textrm{proc} $\Finalize(#1)$}}

\newcommand{\fprff}{\algfont{f}}

\newcommand{\gamefont}[1]{\mathrm{#1}}

\newcommand{\getsr}{{{\,\leftarrow{\hspace*{-3pt}\raisebox{.75pt}{$\scriptscriptstyle\$$}}\,}}}
\newcommand{\bits}{\{0,1\}}
\newcommand{\bit}{\{0,1\}}
\newcommand{\Ex}{\mathbb{E}}
\newcommand{\eqdef}{\stackrel{def}{=}}
\newcommand{\To}{\rightarrow}
\newcommand{\e}{\epsilon}
\newcommand{\R}{\mathbb{R}}
\newcommand{\N}{\mathbb{N}}
\newcommand{\Z}{\mathbb{Z}}
\newcommand{\ie}{\emph{i.e.,} }
\newcommand{\eg}{\emph{e.g.,} }
\newcommand{\VAR}{\mathrm{VAR}}
\newcommand{\Prob}[1]{{\Pr\left[\,{#1}\,\right]}}
\newcommand{\probb}[2]{{\Pr}_{#1}\left[\,{#2}\,\right]}
\newcommand{\probbs}[3]{{\Pr}^{#1}_{#2}\left[\,{#3}\,\right]}
\newcommand{\probbp}[2]{{\Pr}_{#1}'\left[\,{#2}\,\right]}
\newcommand{\Probb}[2]{{{\Pr}_{#1}\left[\,{#2}\,\right]}}
\newcommand{\Probbp}[2]{{{\Pr}_{#1}'\left[\,{#2}\,\right]}}
\newcommand{\condProb}[2]{{\Pr}\left[\,#1\,|\,#2\,\right]}
\newcommand{\CondProb}[2]{{\Pr}\left[\: #1\:\left|\right.\:#2\:\right]}
\newcommand{\CondProbb}[3]{{\Pr}_{#1}\left[\: #2\:\left|\right.\:#3\:\right]}
\newcommand{\expect}[1]{{{\mathbf{E}}[{#1}]}}
\newcommand{\expectt}[1]{{{\mathbf{E}}\left[{#1}\right]}}
\newcommand{\qquadd}{{\quad}}
\newcommand{\next}{ \: ; \: }
\newcommand{\group}{\mathbb{G}}
\newcommand{\true}{\mathsf{true}}
\newcommand{\false}{\mathsf{false}}
\newcommand{\AND}{\;\wedge\;}
\newcommand{\bad}{\mathsf{bad}}
\newcommand{\schemefont}[1]{{\mathsf{#1}}}
\newcommand{\algfont}[1]{{\mathsf{#1}}}
\newcommand{\Adv}{\mathbf{Adv}}
\newcommand{\AdvGen}[3]{\mathbf{Adv}^{\mathrm{#1}}_{#2}(#3)}
\newcommand{\Exp}{\mathbf{Exp}}
\newcommand{\LR}{\mathsf{LR}}

\newcommand{\calA}{{\cal A}}
\newcommand{\calB}{{\cal B}}
\newcommand{\calC}{{\cal C}}
\newcommand{\calF}{{\cal F}}
\newcommand{\calK}{{\cal K}}
\newcommand{\calM}{{\cal M}}
\newcommand{\calO}{{\cal O}}
\newcommand{\calR}{{\cal R}}
\newcommand{\calH}{{\cal H}}
\newcommand{\calG}{{\cal G}}
\newcommand{\calD}{{\cal D}}
\newcommand{\calE}{{\cal E}}
\newcommand{\calP}{{\cal P}}
\newcommand{\calS}{{\cal S}}
\newcommand{\calT}{{\cal T}}
\newcommand{\calU}{{\cal U}}
\newcommand{\calX}{{\cal X}}
\newcommand{\calV}{{\cal V}}
\newcommand{\calW}{{\cal W}}
\newcommand{\calY}{{\cal Y}}

\newcommand{\Keys}{\mathsf{Keys}}
\newcommand{\Dom}{\mathsf{Dom}}
\newcommand{\Rng}{\mathsf{Rng}}

\newcommand{\main}{\procfont{Game}}
\newcommand{\gamesfontsize}{\small}
\renewcommand{\captionfont}{}
\newcommand{\indd}{\hspace*{12pt}}

\newcommand{\mpage}[2]{\begin{minipage}{#1\textwidth} #2 \end{minipage}}
\newcommand{\fpage}[2]{\framebox{\begin{minipage}[t]{#1\textwidth}\setstretch{1.2}\gamesfontsize  #2 \end{minipage}}}
\newcommand{\hfpages}[3]{\hfpagess{#1}{#1}{#2}{#3}}
\newcommand{\hfpagess}[4]{
        \begin{tabular}{c@{\hspace*{.5em}}c}
        \framebox{\begin{minipage}[t]{#1\textwidth}\setstretch{1.2}\gamesfontsize #3 \end{minipage}}
        &
        \framebox{\begin{minipage}[t]{#2\textwidth}\setstretch{1.2}\gamesfontsize #4 \end{minipage}}
        \end{tabular}
    }

\newcommand{\gameName}[1]{\underline{$\main$ #1} \\[2pt]}
\newcommand{\procName}[1]{\underline{#1} \\[2pt]}
\newcommand{\bea}{\begin{eqnarray*}}
\newcommand{\eea}{\end{eqnarray*}}

\newtheorem{thm}{Theorem}[section]
\newtheorem{lem}[thm]{Lemma}
\newtheorem{cor}[thm]{Corollary}
\newtheorem{clm}[thm]{Claim}
\newtheorem{defn}[thm]{Definition}
\newtheorem{rem}[thm]{Remark}

%%%%%%%  Author Notes %%%%%%%
\ifnum\shownotes=1
\newcommand{\authnote}[2]{{ $\ll$\textsf{\footnotesize #2}$\gg$}}
\else
\newcommand{\authnote}[2]{}
\fi
%%%%%%%%%%%%%%%%%%%%%%%%%%%%%%%%%


% end of macros
%%%%%%%%%%%%%%%%%%%%%%%%%%%%%%%%%%%%%%%%%%%%%%%%%%%%%%%%%%%%%%


\title{CS 485: Homework 2}
\date{}

\usepackage{lastpage}
\pagestyle{fancy}
\lhead{\footnotesize \parbox{11cm}{COMPSCI 485, UMass Amherst, Spring 2026} }
\renewcommand{\headheight}{24pt}

\begin{document}

\maketitle
\thispagestyle{fancy}

\noindent \subsection*{2 Phase 2: Annotations analysis:}
\noindent \subsubsection*{2.1 Annotator Feedback} 

In our discussions, this tagging system is reasonable; Many tweets are easy to categorize.\\

\noindent The greatest challenge arises at the boundary between PERS and OTHER. In practice, many tweets express emotions without explicitly identifying a subject or topic, making it difficult to determine whether they should be classified as personal expressions or as unclassifiable. \\

\noindent We also found that a lack of context is a major source of disagreement. Many tweets use nicknames something else, which only make sense if the annotator recognizes them. This particularly affects the distinction between ENT, SPORT, and PERS.

\subsubsection*{2.2 Inter-annotator agreement}
\begin{table}[htbp]
\centering
\begin{tabular}{lccccc}
\hline
      & PERS & ENT & SPORT & SOC & OTHER \\
\hline
PERS  & 310  & 18  & 3     & 8   & 18    \\
ENT   & 47   & 87  & 5     & 2   & 5     \\
SPORT & 5    & 3   & 34    & 1   & 2     \\
SOC   & 13   & 1   & 0     & 14  & 1     \\
OTHER & 74   & 13  & 3     & 2   & 81    \\
\hline
\end{tabular}
\caption{Confusion matrix}
\label{tab:confusion-matrix}
\end{table}

\begin{itemize}
        \item Alex vs Basim
                \begin{itemize}
                        \item $p_e = 0.3287$
                        \item $p_o = 0.716$
                        \item $k = 0.5770$
                \end{itemize}
        \item Alex vs ZiChen
                \begin{itemize}
                        \item $p_e = 0.3527$
                        \item $p_o = 0.684$
                        \item $k = 0.5118$
                \end{itemize}           
        \item Basim vs ZiChen
                \begin{itemize}
                        \item $p_e = 0.3849$
                        \item $p_o = 0.704$
                        \item $k = 0.5187$
                \end{itemize}   
        \item Alex vs Basim vs ZiChen (Fleiss’ kappa)
                \begin{itemize}
                        \item $p_e = 0.3605$
                        \item $p_o = 0.7013$
                        \item $k = 0.5330$
                \end{itemize}   
\end{itemize}              


\subsubsection*{2.3 Aggregate to final dataset}
\noindent If all three annotators agree, that label is adopted; if two out of the three select the same label, the majority vote remains; if the three annotators assign different labels, I will make a decision on the label.\\

\noindent Guidelines: First, if a tweet primarily expresses an emotion but does not address a specific topic, label it as PERS. Second, if the text is too short or too vague, label it as OTHER.

\subsubsection*{3.1 Logistic regression learning}

\noindent \textbf{3.1.1}
\[\frac{\partial l}{\partial \beta} = (1 - g(z))x_1 = 0\]

\noindent No update
\\

\noindent \textbf{3.1.2}
\[\frac{\partial l}{\partial \beta} = (1 - g(z))x_1 = (1-0.6)\times 1 = 0.4\]
\[\beta_1^{new} = \beta_1^{old} + 0.4\]

\noindent Weight increases 0.4\\

\noindent \textbf{3.1.3}
\[\frac{\partial l}{\partial \beta} = (1 - g(z))x_1 = (1-0.999)\times 1 = 0.001\]
\[\beta_1^{new} = \beta_1^{old} + 0.001\]

\noindent Weights increases 0.001\\

\noindent Intuition:
As the loss decreases, there is no longer a need to make significant adjustments to the parameters.\\

\noindent Meaning:
The gradient decreases when the model is on the right track. However, one must also be cautious of vanishing gradients.

\subsubsection*{3.2 Precision and recall: class distribution}

\noindent Accuracy increases; Precision remains unchanged; Recall remains unchanged.\\

\noindent Accuracy is highly influenced by class distribution.
Precision and recall focus more on positive class.

\subsection*{4 Phase 2 continued: NLP experiments}

\subsubsection*{4.1 BOW LogReg}

\subsubsection*{4.1.1 Features and preprocessing}
\begin{itemize}
        \item Features
                \begin{itemize}
\item Baseline: binary bag of words with \texttt{CountVectorizer(binary=True)}
\item Improved Model: TF-IDF with unigram and bigram candidates
                \end{itemize}
        \item Pre-Processing
                \begin{itemize}
                        \item Normalized line breaks and lowercased all text
\item Removed URLs
\item Collapsed repeated whitespace into a single space
                \end{itemize}
        
        \item Normalization
                \begin{itemize}
                        \item Improved model used L2-normalized TF-IDF
\item Baseline do not use additional feature normalization.
                \end{itemize}
        \item Hyperparameter tuning and model selection
                \begin{itemize}
                        \item GridSearchCV and macro-F1
                        \item min\_d, ngram\_range, C, class\_weight
                \end{itemize}
\end{itemize}


\subsubsection*{4.1.2 Results}
\\noindent \textbf{Overall Accuracy:} 0.4560

\begin{table}[h]
\centering
\begin{tabular}{lcc}
\hline
\textbf{Class} & \textbf{Precision} & \textbf{Recall} \\
\hline
ENT   & 0.20 & 0.12 \\
OTHER & 0.31 & 0.91 \\
PERS  & 0.73 & 0.50 \\
SOC   & 0.00 & 0.00 \\
SPORT & 0.00 & 0.00 \\
\hline
\end{tabular}
\end{table}

\newpage
\subsubsection*{4.1.3 Model insight: parameters}
\begin{table}[h]
\centering
\begin{tabular}{p{1.8cm} p{10.5cm}}
\hline
\textbf{Class} & \textbf{Top-10 highest weighted words} \\
\hline
ENT & thi, vs, song, read, want, miss, wa, pleas, dad, na \\
OTHER & fuck, lt, inspir, deserv, congrat, lovelyz, mother, apt, bean, intelectu \\
PERS & don, god, life, say, good, happi, lmao, thing, think, care \\
SOC & straight, sure, exist, actual, peopl, apm, sag, union, repres, npr \\
SPORT & tho, wa, piss, belichick, colleg, footbal, month, problem, stanton, lost \\
\hline
\end{tabular}
\end{table}

\noindent SPORT is related to sports, such as “Belichick,” “football,” and “Stanton”;

\noindent SOC is related to society and politics, such as “union,” “represent,” and “NPR”;

\noindent PERS is related to emotions, such as “life,” “good,” “happy,” and “care.”

\noindent However, some of the higher-weighted words are too general (“thi,” “wa,” “don”), the model has been affected by insufficient data.

\subsubsection*{4.1.4 Model insight: worst mistakes}

\begin{table}[h]
\centering
\begin{tabular}{p{8.5cm} c c c}
\hline
\textbf{Text} & \textbf{True} & \textbf{Pred} & \textbf{Confidence} \\
\hline
the church is homophob the in game pope is fuck gay my man & PERS & OTHER & 0.3201 \\
i swear even with thi fuck up wrist i need that fade from a few bitch nigga i ll prove em a bitch frfr & PERS & OTHER & 0.3183 \\
im not a religi person but i do think sometim god made you for me mariann shut the fuck up sob & PERS & OTHER & 0.2935 \\
\hline
\end{tabular}
\end{table}

\noindent In these three samples, the true labels are all PERS, but the model predicted them all as OTHER. I think the reason is that these examples contain a lot of fuck words, and the model has learned to classify these words as belonging to the OTHER.

\noindent One possible way to improve this is to use a cleaner preprocessing method.


\subsection*{4.2 Improved model}

\begin{table}[h]
\centering
\begin{tabular}{lcc}
\hline
\textbf{Model} & \textbf{Accuracy} & \textbf{Macro-F1} \\
\hline
Baseline: Binary BOW & 0.4560 & 0.2413 \\
Improved: TF-IDF & 0.3840 & 0.2125 \\
\hline
\end{tabular}
\end{table}

\noindent The baseline model achieves higher accuracy (0.4560 vs 0.3840) and a higher macro-F1 score (0.2413 vs 0.2125). TF-IDF is not a very suitable method.

\noindent I think the reason is that the dataset is small and the class distribution is severely imbalanced.


\newpage
\subsection*{4.3 Extra Credit: Bag of embeddings LogReg}
\begin{table}[h]
\centering
\begin{tabular}{lcc}
\hline
\textbf{Model} & \textbf{Accuracy} & \textbf{Macro-F1} \\
\hline
Baseline: Binary BOW & 0.4560 & 0.2413 \\
Improved: TF-IDF & 0.3840 & 0.2125 \\
Glove 25d &  0.5360   &  0.3550\\
GloVe 50d & 0.5360 & 0.3357 \\
GloVe 100d &0.5520   & 0.2903\\
Glove 200d &0.5280   &  0.3292\\
\hline
\end{tabular}
\end{table}

\end{document}
